% Vérifié : citations ramenées à la lecture p2/p1 (Donse, harasse, tete, Double nef, Dû me contenter, Solde), « courte nuit » et Marguerite supprimés (absents ; déjà dans l'introduction générale), pluie du 10 août, logement du 7 août, Maxse, Erquinghem non attesté pour la brigade, ballons du 29 juillet, distance Wytschaete, probabilités atténuées (11th R. Fus., flamand, quinzaine), fin de chapitre corrigée (le chapitre II s'ouvre le 13 août ; anniversaire retiré, embarquement du 25 août f. 14 r°).
% Introduction du chapitre I : « De la Flandre intérieure à la vallée de la Lys » (25 juillet – 12 août 1916 ; clichés 001–018, ff. 1 r° – 9 v°)

Lorsque s'ouvre le carnet, Hippolyte Combes est l'interprète de la 152nd Company de l'\textit{Army Service Corps}, la compagnie hippomobile du train de la 18th (Eastern) Division qui ravitaille la 54th Infantry Brigade («~le ventre d'1 brigade~», selon le titre du chapitre~\textsc{vii} de ses \textit{Feuillets de Route}, annexe, p.~\pageref{annexe:doc6}). La division vient de quitter la Somme, où elle a combattu à Montauban le 1\textsuperscript{er}~juillet 1916, puis au bois des Trônes ; transportée en Flandre, elle passe à la Second Army du général Plumer\footnote{Bedfordshire Archives Service (Bedford), X550/8/1/1607, 1607b et 1608 : \textit{War Diary}, 7th (Service) Battalion Bedfordshire Regiment (54th Infantry Brigade), juillet et août 1916 (ci-après : JMO du 7th Bedfordshire), 23 et 31~juillet 1916 ; G.~H.~F. Nichols, \textit{The 18th Division in the Great War}, Édimbourg et Londres, Blackwood, 1922, p.~51--74.}. Les dix-neuf journées de ce chapitre suivent la brigade de son repos de Flandre intérieure, autour de Lynde, aux environs de Bailleul, puis à La Croix-du-Bac, au bord de la Lys, à quelques kilomètres du front d'Armentières. Pour l'interprète, c'est le service ordinaire d'une compagnie de train au repos puis en secteur : chercher des logements, faire couper le fourrage vert, suivre le convoi. Le ton du carnet est donné dès la première page, dans la notation du ciel de Flandre : «~Mélancolie éparse à travers atmosphère~» (f.~1~r\textsuperscript{o}).

\subsection*{Le carnet}

Les trois premières journées se passent au cantonnement de Lynde, que la mention de sa logeuse, «~M\textsuperscript{me} Donse, de Lind~», permet d'identifier (f.~2~v\textsuperscript{o})\footnote{Le nom de la logeuse, lu aussi «~Donze~», est de lecture incertaine ; celui du village où se fait le rassemblement du 25~juillet, lu «~Sercus~», l'est aussi.}. Le 25~juillet, levé à 8~heures, Combes va à pied avec le «~fantastique lieut\textsuperscript{t} Curtis~» au rassemblement de l'unité, devant une église entourée de son cimetière ; l'arrivée de «~36~off. du R.~Fus\textsuperscript{rs}~» l'oblige à déménager dans une pièce de derrière ; il visite l'église du village, dont il remarque les «~2~rangs de statues~» placés «~sous les orgues~», et reçoit le soir les dames du château, qu'il conduit au mess de ses officiers (f.~1~r\textsuperscript{o}). Le lendemain, de 9~heures à 2~heures et demie, il parcourt le village et ses environs avec le lieutenant Michel, \textit{billeting officer}, pour trouver les trente-six chambres d'officiers, et se trouve libre vers 2~heures, «~harasse de fatigue et mourant de faim~» ; il apprend «~avec regret~» un nouveau changement de cantonnement pour le 28 (f.~1~v\textsuperscript{o}). Le 27, il visite à pied les deux châteaux de Renescure et son église, dont la «~Double nef~» le retient, gagne Hazebrouck à cheval l'après-midi et note, parmi ses achats, des cartes postales d'Ypres, un cahier cartonné et une carte d'état-major ; ses comptes portent aussi une «~Solde : 84~fr.~» (f.~2~r\textsuperscript{o}).

Le 28~juillet est une «~Journée de mouv\textsuperscript{t}~». Parti à 8~heures et demie, Combes chevauche «~en tete du convoi~» par des routes sablonneuses bordées d'avoine et de blé, contourne Hazebrouck par les faubourgs et atteint vers 4~heures et demie une ferme isolée dans une grande prairie, où, faute de chambre, il a «~Dû me contenter d'1~grenier~» (f.~2~v\textsuperscript{o}). Il y reste huit jours, près de Saint-Jans-Cappel. «~Ce fou de Curtis~» l'emmène dès le 29 chercher du fourrage à Saint-Jans-Cappel et au hameau du Mont-Noir ; l'après-midi, en plein soleil, il fait «~faucher \& peser fourrage~» dans un champ ; il s'entretient longuement, avec un interlocuteur que la lecture ne permet pas d'identifier, de «~la dernière attaque~» de la division et décrit la plaine des houblonnières, les ballons captifs à l'horizon, le «~bombardement lointain~» et les avions anglais (f.~3~r\textsuperscript{o}). Le dimanche 30, après une nuit troublée par les rats, il traverse les houblonnières jusqu'à l'église «~archi-comble~» du village, hommes à droite et femmes à gauche, puis lit assis sur un petit pont ; le 31, il va en automobile avec ses officiers au \textit{refilling point} et à Bailleul, et doit l'après-midi conduire un chariot «~to the fields for green forage~» (f.~3~v\textsuperscript{o}). Les jours suivants, souvent libres, sont consacrés à Bailleul, à une demi-heure de marche : la place, que Combes compare à celles d'Hazebrouck et de Béthune (souvenir de son premier service en Artois auprès du \textit{1\textsuperscript{st} Battalion, Irish Guards} en mai--juillet 1915, \textit{Feuillets de Route}, chap.~\textsc{v}--\textsc{vi}, p.~\pageref{annexe:doc6}), l'église Saint-Vaast, l'hôtel de ville, les belles maisons des rues d'Ypres et de Lille, et les dentellières qui travaillent devant leurs portes, dont il se dit «~émerveillé~» (ff.~4~r\textsuperscript{o}--4~v\textsuperscript{o}). Le 4~août, il se rend à cheval, par Bailleul et Steenwerck, à La Croix-du-Bac, où se trouve un «~bureau mission f\textsuperscript{se}~» : il y touche 80~F~30, «~y compris 8\textsuperscript{f}~» de «~supplément~» de juillet (f.~5~r\textsuperscript{o}). Le soir, dans les pâtures de la ferme, il note une «~Impression solitude \& d'abandon~», et entend de son grenier la famille du fermier dire en commun la prière du soir (f.~5~v\textsuperscript{o}).

Le 5~août, le convoi part à 2~heures pour La Croix-du-Bac ; les \textit{horse-lines} s'installent à l'entrée du village, «~du côté de Steenwerck~», et Combes cherche des chambres pour les officiers (f.~6~r\textsuperscript{o}). Le verso de ce feuillet n'a pas été photographié. Le 6, il accompagne Curtis au \textit{refilling point}, «~sur rive canal \& près d'écluse~», visite seul l'église de Bac-Saint-Maur, «~à côté usine~», et écoute l'après-midi, avec un autre interprète, Maurin, une musique militaire anglaise (f.~7~r\textsuperscript{o}). Le 7, revenant d'une matinée de recherche de billets et de \textit{messroom} avec «~cet idiot de C\textsuperscript{tis}~», il trouve la compagnie partie à l'autre bout du village --- «~Damn~!~» --- et s'installe dans une maison «~au dessus Postoffice~», route de Bac-Saint-Maur, avec une grande chambre au premier étage (f.~7~v\textsuperscript{o}). Suivent deux journées de canicule passées dans cette chambre, entre la poussière des routes et les moustiques (f.~8~r\textsuperscript{o}). Le 10, après une matinée de pluie fine, il va à pied jusqu'à Sailly-sur-la-Lys, dont l'église, «~brulée par les Boches~», n'a plus que ses murs et son clocher, et cause longuement à l'écluse avec l'éclusier, notamment du commandant de la division, Maxse ; une mention isolée, «~Les Australiens~», ferme l'entrée (f.~8~v\textsuperscript{o}). Le 11~août, de 2 à 3~heures de l'après-midi, «~40 ou 50~obus~» tombent autour du cantonnement : un soldat est atteint dans la «~Maison du coin, à 50~mètres~», la «~g\textsuperscript{de} blanchisserie Barbry~» subit «~b\textsuperscript{p} de dégâts~», et Combes, descendu à la cave, note «~Enervement et malaise~» et une «~Insécurité perpétuelle~» (f.~9~r\textsuperscript{o}). Le chapitre s'achève le 12 sur une promenade le long de la Lys et un «~repos délicieux~» dans un champ de luzerne (f.~9~v\textsuperscript{o}).

Les comptes du soir dessinent la vie matérielle de l'interprète avant son admission au mess des officiers, en septembre (f.~21~v\textsuperscript{o}) : il paie à ses logeuses la cuisine, les œufs, le lait et la bière, laisse des pourboires au chef cuisinier Hancock (5~F, f.~3~v\textsuperscript{o}) et à son \textit{groom} Connelly (f.~6~r\textsuperscript{o}), et achète des cartes postales presque chaque jour.

\subsection*{Le contexte}

La division arrive en Flandre éprouvée : du 14 au 21~juillet, au bois des Trônes et au bois Delville, elle a perdu 129~officiers et 2~597~hommes\footnote{W.~Miles (comp.), \textit{Military Operations, France and Belgium, 1916}, vol.~II, Londres, Macmillan, 1938, p.~107, n.~3.}. Débarquée à Arques le 23~juillet, la 54th Brigade cantonne autour de Wallon-Cappel, à quatre kilomètres de Lynde\footnote{Distance à vol d'oiseau (4,2~km) entre les centres des deux communes, mesurée par l'éditeur.} ; le 26, Maxse y félicite le 7th Bedfordshire de la prise de la Pommiers Redoubt, le 1\textsuperscript{er}~juillet ; le 28, la brigade marche jusqu'au «~Steam Mill W of Bailleul~» par Hazebrouck, «~a very hot day which tired the men severely~» ; elle s'y entraîne jusqu'au 4~août, puis descend vers la Lys et relève, dans la nuit du 6 au 7, le 1st Wellington, de la 1st New Zealand Brigade, à cheval sur la route d'Armentières à Lille\footnote{JMO du 7th Bedfordshire, 23--31~juillet et 1\textsuperscript{er}--7~août 1916.}. Le déplacement de la compagnie de Combes épouse donc celui de la brigade, et le journal du bataillon confirme jusqu'à la météorologie du carnet : «~Fine hot day~» du 29~juillet au 4~août\footnote{\textit{Ibid.}, 29~juillet -- 4~août 1916.}. Le 7, la compagnie change de cantonnement le matin même où la brigade achève la relève (f.~7~v\textsuperscript{o}). Les «~36~off.~» des Royal Fusiliers sont vraisemblablement ceux du 11th Royal Fusiliers, l'un des quatre bataillons de la brigade\footnote{Nichols, \textit{op.~cit.}, p.~109.}.

Le secteur où descend la division est réputé calme, mais il sort de l'attaque de Fromelles, les 19 et 20~juillet, où la 5th Australian Division a perdu 5~533~hommes\footnote{Miles, \textit{op.~cit.}, p.~119--121 et 133, n.~2 ; le communiqué britannique ne parla que d'«~important raids~» (C.~E.~W. Bean, \textit{The Australian Imperial Force in France, 1916}, 12\textsuperscript{e}~éd., Sydney, Angus \& Robertson, 1941, p.~446).}. Cette division reste en ligne près d'Armentières jusqu'à la mi-octobre ; l'une de ses brigades cantonne à Sailly, et, pour elle, «~the Railhead for supplies was Bac St.~Maur~», d'où le train portait vivres et fourrage aux \textit{refilling points} des brigades\footnote{A.~D. Ellis, \textit{The Story of the Fifth Australian Division}, Londres, Hodder \& Stoughton, [1920], p.~74--76, 117 et 119--121. Que la 18th Division ait eu la même tête de ligne n'est pas établi.} : ce sont les «~Australiens~» et le \textit{refilling point} «~sur rive canal~» du carnet. L'église de Sailly avait brûlé dans les combats d'octobre 1914, et le village resta ensuite à peu près intact jusqu'en avril 1918\footnote{Commonwealth War Graves Commission, historique du cimetière «~Sailly-sur-la-Lys Churchyard~».}. L'arrière n'est pas épargné pour autant : le 9~août, une cinquantaine d'obus lourds tombent sur Armentières, et le 10th Essex note que le quartier général de la division, «~a few miles back~», reçoit «~some attention from enemy artillery~»\footnote{JMO du 7th Bedfordshire, 9~août 1916 ; T.~M. Banks et R.~A. Chell, \textit{With the 10th Essex in France}, 2\textsuperscript{e}~éd., Londres, Gay \& Hancock, 1924, p.~142. Ce quartier général n'est pas localisé ; le rapprochement avec le bombardement du 11~août à La Croix-du-Bac reste une hypothèse.}. Les ballons captifs et le «~bombardement lointain~» notés le 29~juillet près de Saint-Jans-Cappel viennent vraisemblablement du front de Wytschaete et de Messines, à une dizaine de kilomètres à l'est, que les Allemands tiendront jusqu'au 7~juin 1917\footnote{CWGC, historiques des cimetières «~Mont Noir Military Cemetery~», «~Wytschaete Military Cemetery~» et «~Messines Ridge British Cemetery~».}.

Le pays que décrit Combes est celui de la géographie de son temps. Au nord de la voie ferrée d'Hazebrouck à Bailleul, on parle encore flamand : Lynde et Sercus restent «~en majorité germaniques~», Bailleul est «~la citadelle de la langue tudesque en France~», et la prière du soir entendue à Saint-Jans-Cappel était vraisemblablement dite en flamand, la fin de la phrase étant illisible (f.~5~v\textsuperscript{o}) ; les houblonnières sont propres au pied des monts de Flandre, et les dentellières de Bailleul forment la seule industrie dentellière de la Flandre française, «~une centaine d'ouvrières~»\footnote{R.~Blanchard, \textit{La Flandre. Étude géographique de la plaine flamande en France, Belgique et Hollande}, Paris, Armand Colin, 1906, p.~359--360, 382 et 482--483.}. Les «~2~rangs de statues~» de l'église du 25~juillet correspondent vraisemblablement au jubé du \textsc{xv}\textsuperscript{e}~siècle de Saint-Vaast de Lynde, reporté en 1695 au fond de la nef, «~sous les orgues~»\footnote{Ministère de la Culture, base Palissy, notice PM59000995.}. Bailleul, grand centre ferroviaire et hospitalier britannique, sera rasée en 1918, avec l'église Saint-Vaast que visite Combes ; l'église de Steenwerck disparaîtra la même année\footnote{CWGC, historiques des cimetières «~Bailleul Communal Cemetery Extension~» et «~Le Grand Beaumart British Cemetery~» (Steenwerck) ; base Mérimée, notices PA59000232 (Bailleul) et PA59000231 (Steenwerck) ; Églises ouvertes, notice «~Église Saint-Vaast, Bailleul~».}.

Le service décrit par le carnet est celui que les témoins attribuent à l'interprète d'unité. Le cantonnement est «~une des principales besognes de l'interprète~» : un officier, l'interprète et deux caporaux précèdent la troupe, et l'habitant reçoit 1~F par nuit et par lit d'officier\footnote{Fernand-Laurent, \textit{Chez nos alliés britanniques. Notes et souvenirs d'un interprète}, Paris, Boivin, 1917, p.~225--234 ; le \textit{groom}, soldat chargé du cheval d'un officier monté, est décrit p.~64.}. Le fourrage que Combes fait faucher et peser est le \textit{green forage} réglementaire, «~rye grass, lucerne, sainfoin, clover~», qui complète la ration de foin des chevaux\footnote{\textit{Field Service Pocket Book, 1914}, Londres, HMSO, 1914, chap.~6, «~Forage~», §~20.}. La Mission militaire française semble le payer par quinzaine : les 84~F du 27~juillet font seize jours à 5~F~25, tandis que le «~supplément~» de 8~F du 4~août ne se rattache à aucun texte connu\footnote{Calcul de l'éditeur sur les dix-neuf versements du carnet ; le taux journalier de 5~F~25 n'est attesté par aucun texte réglementaire retrouvé.}. Les prix payés aux fermières sont des prix de producteur : le litre de lait à 25~c vaut la moitié du maximum fixé à Paris en février 1917, le litre de bière à 30~c reste au-dessous du tarif d'un économat de chemins de fer en juillet 1916 (44~c)\footnote{L.~March, \textit{Mouvement des prix et des salaires pendant la guerre}, Paris, PUF, et New Haven, Yale University Press, 1925, p.~182--183.}. Ni le lieutenant Curtis, ni Maurin, ni Connelly ne sont identifiés : les dossiers d'officiers des Archives nationales britanniques comptent onze Curtis dans l'\textit{Army Service Corps}, et le carnet ne donne ni prénom ni initiale\footnote{The National Archives (Kew), catalogue Discovery, série WO~339 (dossiers d'officiers), recherche «~Curtis~» et «~Army Service Corps~».}.

Le «~nouveau départ dimanche~» annoncé le 11~août (f.~9~r\textsuperscript{o}) n'eut pas lieu : la compagnie reste près de la Lys jusqu'au 25~août, jour où elle embarque à Bailleul (f.~14~r\textsuperscript{o}). Le 14~août, George~V visite l'école de grenadiers de la 5th Australian Division à Sailly, à deux kilomètres et demi de La Croix-du-Bac, sans que le carnet en dise rien\footnote{Ellis, \textit{op.~cit.}, p.~120--121.}. Viendront ensuite l'arrivée des Highlanders de la 51st Division à Armentières, les 15 et 16~août\footnote{F.~W. Bewsher, \textit{The History of the 51st (Highland) Division, 1914--1918}, Édimbourg et Londres, Blackwood, 1921, p.~87--88.}, puis le transport de la brigade vers le Ternois et, de là, vers la Somme.
